\documentclass[11pt]{article}
\usepackage[margin=1in]{geometry}
\usepackage{amsmath,amssymb}
\usepackage{enumitem}
\usepackage{bookmark}
\hypersetup{
  pdftitle={PSET 3: Introduction to Matrices and Vectors},
  hidelinks}

\title{PSET 3: Introduction to Matrices and Vectors}
\author{}
\date{\vspace{-2.5em}}

% Problems are numbered 1-9 across the document.
\newcounter{problem}
\newcommand{\problem}[1]{\stepcounter{problem}\section*{Problem \theproblem: #1}}

\begin{document}
\maketitle

% Shared assignment questions (header block).
% Include at the top of any problem set with:  % Shared assignment questions (header block).
% Include at the top of any problem set with:  % Shared assignment questions (header block).
% Include at the top of any problem set with:  \input{assignment-qs}
\section*{Assignment Qs}

\begin{itemize}
\item Name
\item How long did this problem set take you (round to nearest hour)? \quad
  0-1 hour \ 2-3 hours \ 4-5 hours \ 5+ hours
\item How difficult was this problem set? \quad
  very easy \ 1 \ 2 \ 3 \ 4 \ 5 \ very challenging
\end{itemize}

\section*{Assignment Qs}

\begin{itemize}
\item Name
\item How long did this problem set take you (round to nearest hour)? \quad
  0-1 hour \ 2-3 hours \ 4-5 hours \ 5+ hours
\item How difficult was this problem set? \quad
  very easy \ 1 \ 2 \ 3 \ 4 \ 5 \ very challenging
\end{itemize}

\section*{Assignment Qs}

\begin{itemize}
\item Name
\item How long did this problem set take you (round to nearest hour)? \quad
  0-1 hour \ 2-3 hours \ 4-5 hours \ 5+ hours
\item How difficult was this problem set? \quad
  very easy \ 1 \ 2 \ 3 \ 4 \ 5 \ very challenging
\end{itemize}


\problem{Sketch a function}

Sketch the graph of a function (any function you like, no need to
specify a functional form) that is:\footnote{inspired by Grimmer HW3.1}

\begin{enumerate}[label=\alph*.]
\item Continuous on \([0,3]\) and has the following properties: an
  absolute maximum at 0, an absolute minimum at 3, a local maximum at 1
  and a local minimum at 2.
\item Do the same for another function with the following properties: 4
  is a \textbf{critical number} (i.e.\ \(f'(x) = 0\) or \(f'(x)\) is
  undefined), but there is no local minimum and no local maximum.
\end{enumerate}

\problem{Find critical values}

Find the critical values of these functions:\footnote{inspired by
  Grimmer HW3.2}

\begin{enumerate}[label=\alph*.]
\item \(f(x) = 5x^{2/3} - 4x\)
\item \(s(t) = 3t^4 - 4t^3 + 6t^2\)
\item \(f(r) = \dfrac{r}{r^2 + r + 1}\)
\item \(h(x) = x \ln(x)\)
\end{enumerate}

\problem{Find absolute minimum/maximum values}

Find the absolute minimum and absolute maximum values of the functions
on the given interval:\footnote{inspired by Grimmer HW3.3}

\begin{enumerate}[label=\alph*.]
\item \(f(x) = 3x^2 - 12x + 5\), \([0,1]\)
\item \(f(t) = t^2\sqrt{9 - t^2}\), \([-1,4]\)
\item \(s(x) = x - \ln(x)\), \([1/2, 2]\)
\end{enumerate}

\problem{Basic matrix arithmetic}

If
\[
\mathbf{a} = \begin{bmatrix} 2 \\ 2 \end{bmatrix}
\quad \mbox{and} \quad
\mathbf{b} = \begin{bmatrix} 3 \\ 1 \end{bmatrix}
\]
find:\footnote{inspired by Pemberton and Rau 11.1.2}

\begin{enumerate}[label=\alph*.]
\item \(\mathbf{a} - \mathbf{b}\)
\item \(-4 \mathbf{b}\)
\item \(3 \mathbf{a} + 3 \mathbf{b}\)
\end{enumerate}

\problem{More complex matrix arithmetic}

Suppose
\[
\mathbf{x} = \begin{bmatrix} 3 \\ 2q \\ 4 \end{bmatrix}
\quad \mbox{and} \quad
\mathbf{y} = \begin{bmatrix} p + 2 \\ -5 \\ 3r \end{bmatrix}
\]
If \(\mathbf{x} = 2 \mathbf{y}\), find \(p, q, r\).\footnote{inspired by
  Pemberton and Rau 11.1.3}

\problem{Check for linear dependence}

Which of the following sets of vectors are linearly
dependent?\footnote{inspired by Pemberton and Rau 11.1.4} In each part,
you can denote each vector as \(\mathbf{a}, \mathbf{b}, \mathbf{c}\)
respectively.

\begin{enumerate}[label=\alph*.]
\item \(\begin{bmatrix} 1 \\ 2 \\ 3 \end{bmatrix}, \begin{bmatrix} 4 \\ 5 \\ 6 \end{bmatrix}, \begin{bmatrix} 7 \\ 8 \\ 9 \end{bmatrix}\)
\item \(\begin{bmatrix} 13 \\ 7 \\ 9 \\ 2 \end{bmatrix}, \begin{bmatrix} 0 \\ 0 \\ 0 \\ 0 \end{bmatrix}, \begin{bmatrix} 3 \\ -2 \\ 5 \\ 8 \end{bmatrix}\)
\item \(\begin{bmatrix} 1 \\ 2 \\ 1 \end{bmatrix}, \begin{bmatrix} 2 \\ -2 \\ -1 \end{bmatrix}, \begin{bmatrix} 2 \\ 1 \\ 3 \end{bmatrix}\)
\end{enumerate}

\problem{Vector length}

Find the length of the following vectors:\footnote{inspired by Simon and
  Blume 10.10}

\begin{enumerate}[label=\alph*.]
\item \((4, 2)\)
\item \((2, 2, 2)\)
\item \((4, 2, 3, 1)\)
\item \((0, 0, 0, 0, 3)\)
\end{enumerate}

\problem{Law of cosines}

The \textbf{law of cosines} states:
\[
\cos(\theta) = \frac{\mathbf{v} \cdot \mathbf{w}}{\|\mathbf{v}\| \: \|\mathbf{w}\|}
\]
where \(\theta\) is the angle from \(\mathbf{w}\) to \(\mathbf{v}\)
measured in radians. Of importance, \(\arccos()\) is the inverse of
\(\cos()\):
\[
\theta = \arccos \left( \frac{\mathbf{v} \cdot \mathbf{w}}{\|\mathbf{v}\| \: \|\mathbf{w}\|} \right)
\]
For each of the following pairs of vectors, calculate the angle between
them. Report your answers in both radians and degrees. To convert
between radians and degrees:\footnote{inspired by Simon and Blume 10.12}
\[
\text{Degrees} = \text{Radians} \times \dfrac{180^{\circ}}{\pi}
\]

\begin{enumerate}[label=\alph*.]
\item \(\mathbf{v} = (1, 0), \quad \mathbf{w} = (2, 2)\)
\item \(\mathbf{v} = (4, 1), \quad \mathbf{w} = (2, -8)\)
\item \(\mathbf{v} = (1, 1, 0), \quad \mathbf{w} = (1, 2, 2)\)
\end{enumerate}

\problem{Matrix algebra}

Using the matrices below, calculate the following. Some may not be
defined; if that is the case, say so.\footnote{inspired by Grimmer HW5.3}
\[
\mathbf{A} = \begin{bmatrix} 3 \\ -2 \\ 9 \end{bmatrix}
\quad
\mathbf{B} = \begin{bmatrix} 8 \\ 0 \\ -1 \end{bmatrix}
\quad
\mathbf{C} = \begin{bmatrix} 7 & -1 & 5 \\ 0 & 2 & -4 \end{bmatrix}
\quad
\mathbf{D} = \begin{bmatrix} 3 & 1 \\ 3 & 4 \\ 3 & -7 \end{bmatrix}
\quad
\mathbf{E} = \begin{bmatrix} 5 & 2 & 3 \\ 1 & 0 & -4 \\ -2 & 1 & -6 \end{bmatrix}
\]
\[
\mathbf{F} = \begin{bmatrix} 4 & 1 & -5 \\ 0 & 7 & 7 \\ 2 & -3 & 0 \end{bmatrix}
\quad
\mathbf{G} = \begin{bmatrix} 2 & -8 & -5 \\ -3 & 7 & -4 \\ 1 & 0 & 3 \\ 1 & 2 & 6 \end{bmatrix}
\quad
\mathbf{K} = \begin{bmatrix} 9 \\ -2 \\ -1 \\ 0 \end{bmatrix}
\quad
\mathbf{L} = \begin{bmatrix} 5 & 0 & 3 & 1 \end{bmatrix}
\quad
\mathbf{M} = \begin{bmatrix} 1 & -1 \\ -1 & 3 \end{bmatrix}
\]

\begin{enumerate}[label=\alph*.]
\item \(\mathbf{A} + \mathbf{B}\)
\item \(-\mathbf{G}\)
\item \(\mathbf{D}'\)
\item \(\mathbf{C} + \mathbf{D}\)
\item \(\mathbf{A}' \mathbf{B}\)
\item \(\mathbf{BC}\)
\item \(\mathbf{FB}\)
\item \(\mathbf{E} - 5\mathbf{I}_3\)
\item \(\mathbf{M}^2\)
\end{enumerate}

% Shared AI and Resources statement.
% Include in any problem set with:  % Shared AI and Resources statement.
% Include in any problem set with:  % Shared AI and Resources statement.
% Include in any problem set with:  \input{ai-statement}
\section*{AI and Resources statement}

Please list (in detail) all resources you used for this assignment. If
you worked with people, list them here as well. It is not enough to say
that you used a resource for help, you need to be specific on the link
and \emph{how} it was helpful. W/R/T gen AI tools (including GPT,
Claude, etc.) you cannot use them to do work on your behalf---you
cannot put in any of the questions, etc. You can ask for help on logic
/ sample problems. If you do use GPT or other AI tools, you need to
provide a link to your chat transcript. Any suspected academic
integrity violations will be immediately reported to the Dean of
Students.

\section*{AI and Resources statement}

Please list (in detail) all resources you used for this assignment. If
you worked with people, list them here as well. It is not enough to say
that you used a resource for help, you need to be specific on the link
and \emph{how} it was helpful. W/R/T gen AI tools (including GPT,
Claude, etc.) you cannot use them to do work on your behalf---you
cannot put in any of the questions, etc. You can ask for help on logic
/ sample problems. If you do use GPT or other AI tools, you need to
provide a link to your chat transcript. Any suspected academic
integrity violations will be immediately reported to the Dean of
Students.

\section*{AI and Resources statement}

Please list (in detail) all resources you used for this assignment. If
you worked with people, list them here as well. It is not enough to say
that you used a resource for help, you need to be specific on the link
and \emph{how} it was helpful. W/R/T gen AI tools (including GPT,
Claude, etc.) you cannot use them to do work on your behalf---you
cannot put in any of the questions, etc. You can ask for help on logic
/ sample problems. If you do use GPT or other AI tools, you need to
provide a link to your chat transcript. Any suspected academic
integrity violations will be immediately reported to the Dean of
Students.


\end{document}
